\documentclass{article}

\usepackage{graphicx} % Required for inserting images

%\usepackage{graphicx} % Required for inserting images


\usepackage{amsmath,amssymb,amsthm}
\usepackage{mathrsfs}
\usepackage{tikz}
\usetikzlibrary{cd}
\usepackage{ upgreek }
\usepackage{ tipa }
\usetikzlibrary{cd}
\usepackage{dsfont}
\usepackage{extpfeil}
\usepackage[final]{todonotes}
\usepackage{xcolor}
\usepackage{bbold}
\usepackage{tikz}
\usetikzlibrary{cd}

\usepackage{hyperref}
\usepackage[notcite,notref,final]{showkeys}

\newcommand{\C}{\mathds{C}}
\newcommand{\cC}{\mathcal{C}}
\newcommand{\cT}{\mathcal{T}}
\newcommand{\cD}{\mathcal{D}}

\newcommand{\adjunction}[4]{\begin{tikzcd}[ampersand replacement=\&] #1 : #2 \arrow[r, shift left] \& #3 : #4 \arrow[l, shift left] \end{tikzcd}}
\newcommand{\Dperf}{D^{\textrm{\textup{perf}}}}
\newcommand{\module}[2][]{#2\text{-}\modop_{#1}}
\newcommand{\prsp}{\mathbb{P}}
\newcommand{\set}[2]{ \left\{ #1\, \middle| \,#2 \right\}}
\newcommand{\setl}[1]{\left\{ #1 \right\}}
\newcommand{\strshf}[1]{\mathscr{O}_{#1}}
\newcommand{\todoinline}[1]{\todo[inline]{#1}}

\newcommand{\vivek}[1]{\textcolor{blue}{\textbf{Vivek}: #1}}
\newcommand{\cA}{\mathscr{A}}
\newcommand{\chA}{\Ch(\cA)}
\newcommand{\chX}[1][X]{\Ch(\QCoh(#1))}

\DeclareMathOperator{\Ch}{Ch}
\DeclareMathOperator{\colim}{colim}
\DeclareMathOperator{\Ho}{Ho}
\DeclareMathOperator{\id}{id}
\DeclareMathOperator{\modop}{mod}
\DeclareMathOperator{\QCoh}{QCoh}


\newtheorem{theorem}{Theorem}[section]
\newtheorem{lemma}[theorem]{Lemma}

\newtheorem*{theorem*}{Theorem}
\newtheorem{proposition}[theorem]{Proposition}
\theoremstyle{definition}
\newtheorem{definition}[theorem]{Definition}
\newtheorem{notation}[theorem]{Notation}
\newtheorem*{notation*}{Notations}
\newtheorem{problem}[theorem]{Problem}
\theoremstyle{remark}
\newtheorem{remark}[theorem]{Remark}
\newtheorem{example}[theorem]{Example}
\newtheorem{corollary}[theorem]{Corollary}
\newtheorem{assumption}[theorem]{Assumptions}

\title{Stable compatibility for model categories}
\author{Kunal Arora \and Vivek Mohan Mallick}


\begin{document}

\maketitle
\begin{abstract}
By introducing the notion of stable compatibility, we provide a framework in the model category theory which addresses the existence of a stable model structure on the category of modules over a monad. Exact version of such monads have been studied in the setting of triangulated categories, separately as separable monads \`a la Balmer in tensor triangular geometry and non-separable monads \`a la Dell'Ambrogio and Sanders in stable homotopy theory. By introducing stable compatibility, we explore their properties together and produce a non-separable monad in Algebraic geometry and more.
\end{abstract}
%\keywords{modules over monads, stable model category, derived categories, spectra, triangulated categories, etale maps, finite and flat maps, tensor triangular geometry}
%\subjclass{14A30, 55P42, 18G80, 18N40}

\tableofcontents

\section{Introduction}
We are interested in the following broad question:
\begin{quote}
Given a stable model category  $\cC$ and a monad  $(M, \mu, \eta)$ acting on $\cC$, under what conditions is $\module{M}$, the category of modules of $M$ in $\cC$ i.e.\ the Eilenberg-Moore category, a stable model category?
\end{quote}
This interests us majorly in the case of Algebraic geometry where $\cC$ is $\Ch(QCoh(X))$, the category of chain complexes over quasi-coherent sheaves and $M$ is the monad generated by the adjunction induced by a map of quasi-compact and quasi-separated schemes $f:V \rightarrow X$.

The derived version of the above question is a result of Balmer (See, \cite{balmer:etalneth} theorem 3.5 for more details) about the monad $M=Rf_{*}f^{*}$ generated by a seperated \'etale morphism $f:V \rightarrow X$ of  quasi-compact and quasi-seperated schemes acting on the unbounded derived category of quasi-coherent sheaves $D(QCoh(X))$ and it states that $\module{M}$  is triangulated and equivalent to $D(QCoh(V))$. 

This opens up the definition of \'etale morphism and hence \'etale-topology in the setting of \textit{tensor-triangular geometry}.

Our original motivation for this paper was to see what other morphisms and topologies from algebraic geometry  admit an extension in  \textit{tensor-triangular geometry}.

The monad associated to the seperated \'etale map is an example of \textit{separable monad}. We now state an example of a \textit{non-separable} monad from stable homotopy theory.

A similar situation arises in stable homotopy theory, when one considers the map say $f: R \rightarrow A$ between highly structured ring spectrum (say \cite{MR1417719} $S$-algebra), this \textit{change of ring} map induces an extension of scalars and restriction adjunction on the level of module categories $f^{*}=A \wedge - : M_{R} \leftrightarrows M_{A}:f_{*} $ where  $(f^{*} \dashv f_{*})$. Let $R$ be the sphere spectrum, denoted as $S$, and let $A$ be as above with $f: S \rightarrow A$ the unit map. The category of modules of the monad $f_{*}f^{*}$ on $M_{S}$ is $M_{A}$ which inherits a right-transferred model structure from $M_{S}$ and is simplicial, stable and cofibrantly generated model category.

Consider the monad formed by the derived adjunction $Lf^{*}: D(S) \leftrightarrows D(A): Rf_{*}$ between $D(S)=SH$ the stable homotopy of spectra and the derived category $D(A)$. Let $A\text{-Mod}_{SH}$ be the category of modules of the monad $Rf_{*}Lf^{*}$ on $D(S)$. Dell'Ambrogio and Sanders \cite[corollary 1.11]{MR3851535} characterizes the existence of triangulated structure on $A\text{-Mod}_{SH}$ in terms of its (exact) equivalence with $D(A)$. 
Example 1.12 (loc.\ cit.) says that Guti\'errez shows that corollay 1.11 holds for  $A=HB$ an Eilenberg Maclane spectrum of a field or subring B of $\mathbb{Q}$. 
    

We note that the case of stable homotopy which provides an example of a non-separable monad says that we get a triangulation on  $A\text{-Mod}_{SH}$ only when it is equivalent to the homotopy category of a stable model category $M_{A}$. This lead us to investigate the question mentioned in the beginning and to provide a unified treatment for both separable and non-separable monads on stable model categories. 


\subsection*{Main results}
Answering the above mentioned question packed in the definition of \textit{stable compatibility} (see definition \ref{def:stabcomp}), which arises when the forgetful functor preserves cofibrant objects under some suitable hypothesis, we have our first main result (see theorem \ref{the:stabcomp} below).
 
\begin{theorem*}
    
Let $\cC$ be a stable cofibrantly generated model category.
  Suppose $M$ is a monad on $\cC$ which satisfies the following.
  \begin{enumerate}
    \item
      The category $\module{M}$ admits a right transferred model structure, which is 
      cofibrantly generated, and such that the induced adjunction $\Sigma_{\module{M}}
      \dashv \Omega_{\module{M}}$ is a Quillen adjunction. 
    \item The monad $M$ satisfies stable compatibility.
  \end{enumerate}
  Then $\module{M}$ is a stable model category. Moreover, if the monad is induced by a
  Quillen adjunction $\adjunction{F}{\cC}{\cD}{G}$, the underlying comparison functor $K$
  is such that $K \circ F=F_{M}$ and $G=G_{M} \circ K$ as in the diagram below   
  \[
    \begin{tikzcd}
      & \cC  \arrow[dr,"F_{M}", shift left=1mm] \arrow[d,"F", shift left=1mm] \\
      & \cD \arrow[dashrightarrow,r,"\exists !K" below] \arrow[u,"G",shift left=1mm]& \module{M} \arrow[ul, "G_{M}",shift left=1mm] 
    \end{tikzcd}
  \]
  and it preserves fibrations and acyclic fibrations.


\end{theorem*}
We use this result and its interaction with sections of a separable monad to obtain our second main result (see theorem \ref{the:mainsepm}).

\begin{theorem*}
    Let $\cC$ be a cofibrantly generated stable model category, and let $M$ be a 
separable monad on $\cC$ such that $M$ satisfies the assumptions of theorem 
\ref{the:stabcomp}. Then $\module{M}$ is a cofibrantly generated stable model 
category. Moreover, a cofiber sequence admits a direct summand inclusion into a suitable
cofiber sequence.
\end{theorem*}
Although this result provides a clarification to \cite[remark 6.1]{balmer:septrcat} but we mention it now that this does not address the question about the existence of an equivalence between $\Ho(\module{M})$ and $\module[\Ho \cC]{(\Ho M)}$ where the former is the homotopy category of $M$-modules and latter is the category of modules over the derived monad in the homotopy category .

This result provides us with a model categorical insight about the structure of data of cofiber sequences (see remark \ref{rem:sepscmnd}), and for an application to the derived category $D(S)$ of an \'etale  $R$-algebra $S$ for a commutative ring $R$, see example \ref{exa:sepalgeb}.

Another strong application of \textit{stable compatibility} is in algebraic geometry where we show that finite and flat morphisms between suitable schemes form a monad whose module category turns out to be a cofibrantly generated stable model category (see \textsection \ref{ssc:finflmor}). In this case, we even have an equivalence between between $\Ho(\module{M})$ and $\module[\Ho \cC]{(\Ho M)}$. This is a first major step outside \'etale morphisms in \textit{tensor triangular geometry}.


We mention other potential applications and conclude introduction in this section-wise overview. 

Section \ref{sec:prereqst} provides the preliminaries on Eilenberg-Moore categories $\module{M}$, simplicial model categories, the pre-triangulated structure in the sense of \cite{MR1650134} and stable model categories. We also mention about the existence of right-transferred model structures on $\module{M}$, and lastly we recall the locally free model structure on $\Ch(QCoh(S))$ of a suitable scheme. 

Section \ref{sec:stablecpt} is our main section on \textit{Stable compatibility} with our main result theorem \ref{the:stabcomp} and an auxiliary lemma \ref{lem:cndforsc}, which helps us in detecting scenarios where \textit{stable compatibility} holds. It forms our main tool in \textsection \ref{ssc:finflmor}. An analogous result, concluding stability of the module category in the $\infty$-category setup, can be deduced from \cite[corollary 5.5 and theorem 5.7]{MR3338539}. 

Section \ref{sec:sepmonad} addresses the interaction of \textit{stable compatibilty} with the definition of a \textit{separable monad}.

Section \ref{sec:applications} mentions other potential applications in stable homotopy (see \textsection \ref{ssc:drumhack}), which comes along with an improved observation to a result of [\cite{MR3937590} Theorem 2.3]. Lastly, we mention that in the dual setup of \textit{stable compatibility}, one can recover the stability of the category of comodules over a comonad in \textsection \ref{ssc:comodules}.



\section{Prerequisites} \label{sec:prereqst}

\subsection{Module categories} \label{ssc:monads}
Throughout this paper, a monad on a category $\cC$ is denoted by a triple $(M,\mu,\eta)$ (or simply $M$, when the context is clear). By $\module[\cC]{M}$ (or simply, $\module{M}$ when there is no ambiguity), we denote the category of $M$-modules, the Eilenberg-Moore category \cite{maclane:catworkmath} associated with $M$. The free-forgetful adjunction is denoted by 
\[
\adjunction{F_M}{\cC}{\module[\cC]{M}}{G_M}.
\]
The following definitions are taken from  \cite{balmer:septrcat}.

\begin{definition}(\cite{balmer:septrcat}, Definition 2.1)\label{def:monad}  Let $\cC$ be a category. A \textit{monad} on $\cC$ is a triple $(M,\mu,\eta)$ where $M: \cC \rightarrow \cC$ is an endofunctor, $\mu:M^{2} \rightarrow M$ (\textit{the multiplication}) and $\eta:Id_{\cC} \rightarrow M$ (\textit{the unit}) are natural transformations such that the following diagrams expressing associativity and two-sided unit,commute:
\begin{equation}
    \begin{tikzcd}
      M^{3} \arrow{r}{M\mu} \arrow[swap]{d}{\mu M} & M^{2} \arrow{d}{\mu} \\
      M^{2} \arrow{r}{\mu} & M
      \end{tikzcd}
     and 
     \begin{tikzcd}
      M \arrow{r}{M\eta} \arrow[swap]{dr}{id} & M^{2} \arrow{d}{\mu} & M \arrow[swap]{l}{\eta M} \arrow{dl}{id} \\
       &M .
     \end{tikzcd}
\end{equation}
When the category $\cC$ is suspended(i.e. an additive category with an autoequivalence $\Sigma:\cC \xrightarrow{\sim}\cC$, see \cite{balmer:septrcat}  Definition 1.1 for more), we say that an additive monad $M$ is \textit{stable} if $M$, $\mu$, $\eta$ commute with suspension: $\Sigma M=M \Sigma , \mu \Sigma=\Sigma \mu , \eta\Sigma=\Sigma\eta.$
\end{definition}
See also the assumptions \ref{asm:stabcomp} below.%\todo{clarify the todo there}
\begin{example}
    (See, \cite{balmer:septrcat} Example 2.3) Let $\cC$ be a monoidal category with tensor $\bigotimes: \cC \times \cC \rightarrow \cC$ and the unit $\mathbb{1}  \in \cC$. A (unital and associative) ring object in $\cC$ is a triple $(A,\mu,\eta)$ where $A$ is an object of $\cC$ and where the \textit{multiplication} $\mu: A \bigotimes A \rightarrow A$ and the unit $\eta: \mathbb{1} \rightarrow A$ are morphisms in $\cC$ satisfying the usual associativity and two-sided unit conditions analogous to Definition 2.1 (1). Let $M:\cC \rightarrow \cC$ be the functor $A \bigotimes-$, with the obvious $\mu: M^{2} \rightarrow M$ and $\eta:Id \rightarrow M.$ Then $M$ is a monad on $\cC$.

\end{example}
For a commutative ring $R$ and for $\cC=R-Mod$ with $\bigotimes=\bigotimes_{R}$, the ring objects in $\cC$ are the usual $R$-algebras. In particular, ring objects in $\mathbb{Z}-Mod$ are ordinary rings. So, monads are generalizations of rings and algebras.


\begin{definition}(See, \cite{balmer:septrcat} Definition 2.4)\label{def:M-mod} Let $M:\cC \rightarrow \cC$ be a monad. The \textit{Eilenberg-Moore category} of (left)$M$-modules $M-mod$ is defined as follows. A left $M$-module is a pair $(x, \lambda)$ where $x$ is an object of $\cC$ (\textit{the underlying object}) and $\lambda:Mx \rightarrow x$ is a morphism(\textit{the left action}) such that the following diagrams both commute:
\[\begin{tikzcd}
M^{2}x \arrow{r}{M\lambda} \arrow[swap]{d}{\mu_{x}} & Mx \arrow{d}{\mu} \\
Mx \arrow{r}{\lambda} & x
\end{tikzcd}
\] and 
\[\begin{tikzcd}
x \arrow{r}{\eta_{x}} \arrow[swap]{dr}{id} & Mx \arrow{d}{\lambda} \\
& x
\end{tikzcd}
\]
A morphism $f: (x,\lambda) \rightarrow (x^{'}, \lambda')$ of left $M$-modules is a morphism $f: x \rightarrow x^{'}$ in $\cC$ which is $M$-linear, i.e. such that $\lambda'\circ M(f)=f \circ \lambda$. The \textit{free module} functor $F_{M}: \cC \rightarrow M-Mod$ defined by $F_{M}(y):=(My,\mu_{y})$. The right adjoint $G_{M}:M-Mod \rightarrow \cC$ just forgets the action. We define the \textit{Klesli category of free $M$-modules} $M-Free_{\mathcal{C}}$ as the full subcategory $F_{M}(\cC)$ of  $M-Mod$. The above functors $F_{M}$ and $G_{M}$ restrict to an adjunction between $\cC$ and $M-Free_{\cC}$: \[\begin{tikzcd}
 & \cC \arrow[dl,"F_{M}"] \arrow[dr,"F_{M}"]  \\
 M-Free_{\cC} \arrow[hookrightarrow, rr,"(fully faithful)"below] \arrow[ur,"G_{M}",shift left=1ex] && M-Mod \arrow[ul, "G_{M}",shift left=1ex] 
\end{tikzcd}
\]
When $\mathcal{C}$ is suspended and $M$ is stable, then both $M-Free_{\mathcal{C}}$ and $M-Mod$ inherit an obvious suspension such that $F_{M}$ and $G_{M}$ commute with that suspension.

\end{definition}
Contrast the \textit{stable monad} with the  assumptions \ref{asm:stabcomp} and Remark 3.9 below.
\begin{example}
    (See, \cite{balmer:septrcat} Example 2.5) When $(A,\mu,\eta)$ is a ring object in a monoidal category $\cC=(\cC, \bigotimes,\mathbb{1})$ as in example 2.2, the above constructions yield the natural categories of left $A$-modules and free left $A$-modules.
\end{example}
\begin{remark}(See, \cite{balmer:septrcat} Remark 2.7)%\todo{Redundant info}
    Given an adjunction $F:\mathcal{C} \leftrightarrows \mathcal{D}: G$, let $\eta: Id_{\mathcal{C}} \rightarrow GF$ be the unit and $\epsilon:FG \rightarrow Id_{\mathcal{D}}$ be the counit of this adjunction. Then $M:=GF$ is a monad on $\mathcal{C}$, with unit $\eta$ and multiplication $\mu:=G\epsilon F:M^{2}=G(FG)F \xrightarrow{G \epsilon F} GF=M$. One says that $M$ is \textit{realized by the adjunction} $(F,G)$. 
    
   
    
    If moreover $\mathcal{C}$ and $\mathcal{D}$ are suspended and $F$ and $G$ commute with suspension, we tacitly assume that the adjunction respects the suspension, meaning that $\eta$ and $\mu$ commute with suspension. In that case, the monad $GF$ is stable.
\end{remark}

Given a monad $M:\cC \rightarrow \cC$ , there are in general many adjunctions realizing $M$. They form a category in which the Klesli construction $M-Free_{\mathcal{C}}$ is initial and the Eilenberg-Moore category $M-Mod$ is final, as stated precisely below.

\begin{proposition}[\cite{balmer:septrcat}, proposition 2.8] \label{pro:kleisliem}
Let $M:\cC \rightarrow \cC$  be a monad  \textit{realized by an adjunction }$F:\mathcal{C} \leftrightarrows \mathcal{D}: G$. Then there are unique functors L and K as follows: \[\begin{tikzcd}
 & \cC \arrow[dl,"F_{M}"] \arrow[dr,"F_{M}"] \arrow[d,"F"] \\
 M-Free_{\cC} \arrow[dashrightarrow,r,"\exists !L"below] \arrow[ur,"G_{M}",shift left=1ex] & D \arrow[dashrightarrow,r,"\exists !K"below] \arrow[u,"G",shift left=1ex]& M-Mod \arrow[ul, "G_{M}",shift left=1ex] 
\end{tikzcd}
\]
such that $L \circ F_{M}=F, G_{M}=G \circ L$
and $G=G_{M} \circ K$. Moreover, $L$ is fully faithful. 
Finally, if $\mathcal{C}$ and $\mathcal{D}$ are suspended, if $M$ is stable and if $(F,G)$ is an adjunction of functors commuting with suspension, then the functors $L$ and $K$ commute with suspension as well.

\end{proposition}
Contrast this with the Theorem \ref{the:stabcomp}below.
\subsection{Model categories}
Throughout this paper our model categories are as defined in \cite{MR1650134} definition 1.1.3 (See for more details). The exciting feature of this definition is that the functorial factorizations are a part of the model structure. 
\begin{definition}\label{def:pointedmodcat}
    (As on page 4, \cite{MR1650134}) We call a model category(or any category with an initial and terminal object) \textit{pointed} if the map from the initial object to terminal object is an isomorphism. Given a model category $\cC$, define $\cC_{*}$ to be the category under the terminal object $*$. That is, an object of $\cC_{*}$ is a map $* \xrightarrow{v} X$ of $\cC$, often written $(X,v)$. We think of $(X,v)$ as an object together with a \textit{basepoint} $v$. A morphism from $(X,v)$ to $(Y,w)$ is a morphism $X \rightarrow Y$ of $\cC$ that takes $v$ to $w$.
\end{definition}
By (\cite{MR1650134}, Proposition 1.1.8) one can induce a model structure on $\cC_{*}$ from $\cC$ via the adjunction forgetful functor $U: C_{*} \rightarrow C$ and it's left adjoint of adding the basepoint i.e. X goes to $X_{+}=X \sqcup *$.
\\Let's recall some definitions needed to define \textit{Cofibrantly generated model categories}
\begin{definition}
    (\cite{MR1650134}, Definition 2.1.3) Suppose $\cC$ is a category with all small colimits, $\cD$ is a collection of morphisms in $\cC$, $A$ is an object of $\cC$ and $\kappa$ is a cardinal. We say that $A$ is $\kappa$-\textit{small relative} to $\cD$ if, for all $\kappa$-filtered ordinals $\lambda$ and all $\lambda-$sequences $X_{0} \rightarrow X_{1} \rightarrow...\rightarrow X_{\beta}\rightarrow ...$ such that each map $X_{\beta} \rightarrow X_{\beta+1}$ is in $\cD$ for $\beta+1 < \lambda$, the map of sets $\colim_{\beta < \lambda} \cC(A,X_{\beta}) \rightarrow \cC(A, \colim_{\beta < \lambda} X_{\beta})$ is an isomorphism. We say that $A$ is \textit{small relative} to $\cD$ if  it is $\kappa-$small relative to $\cD$ for some $\kappa$. We say that $A$ is small if it is small relative to $\cC$ itself.
\end{definition}
From (\cite{MR1650134}, Definition 2.1.7) we state that for $I$ a class of maps in $\cC$, a map is \textit{$I$-injective} if it has the right lifting property with respect to every map in $I$. The class of \textit{$I$-injective} maps is denoted by $I$-inj. Similarly, a map with the left lifting property with respect to the class $I$-inj is called \textit{$I$-cofibration} and the class of such maps is denoted by $I$-cof. 
\begin{definition}
    (\cite{MR1650134}, Definition 2.1.9)\label{def:I-cell} Let $I$ be a set of maps in a category $\cC$ containing all small colimits. A \textit{relative $I$-cell complex} is a transifinite composition of pushouts of elements of $I$ i.e. if $f: A \rightarrow B$ is a relative $I$-cell complex, then there is an ordinal $\lambda$ and a $\lambda-$sequence such that $f$ is the composition of $X$ and such that, for each $\beta$ with $\beta+1 < \lambda$, there is a pushout square
     \[\begin{tikzcd}
        C_{\beta} \arrow{r} \arrow{d}{g_{\beta}}  & X_{\beta} \arrow{d} \\
        D_{\beta} \arrow{r} & X_{\beta+1}
     \end{tikzcd}
     \]  
     such that $g_{\beta} \in I$. We denote the collection of relative $I$-cell complexes by \textit{$I$-cell}.
\end{definition}
We now state (\cite{MR1650134}, Theorem 2.1.14) in the form of a lemma.
\begin{lemma}
    (The small object arguement) Suppose $\cC$ is a category containing all small colimits, and $I$ is a set of maps in $\cC$. Suppose the domains of the maps of $I$ are small relative to $I$-cell. Then there is a functorial factorization $(\gamma,\delta)$ on $\cC$ such that, for all morphisms $f$ in $\cC$, the map $\gamma (f)$ is in $I$-cell and the map $\delta(f)$ is in $I-inj$. 
\end{lemma}
We now recall the definition of a \textit{Cofibrantly generated model category}. 
\begin{definition} \label{def:cofgenemodcat}
    (\cite{MR1650134}, Definition 2.1.17) Suppose $\cC$ is a model category. We say that $\cC$ is \textit{cofibrantly generated} if there are sets $I$ and $J$ of maps such that:
    \begin{enumerate}
        \item The domains of the maps of $I$ are small relative to $I$-cell;
        \item The domains of the maps of $J$ are small relative to $J$-cell;
        \item The class of fibrations is $J-inj$; and 
        \item The class of trivial fibrations is $I-inj$.
    \end{enumerate}
    We refer to $I$ as the set of \textit{generating cofibrations}, and to $J$ as the set of \textit{generating trivial cofibrations}.
\end{definition}
We note that the small object arguement of lemma 2.10 gives the functorial factorizations in a \textit{cofibrantly generated model category}. For more details see (\cite{MR1650134}, Proposition 2.1.18) and also Theorem 2.1.19).

\subsection*{Simplicial model categories}

We now recall the notion of a \emph{pointed simplicial model category}. Let $SSet$ be the category of simplicial sets and $SSet_{*}$ be the pointed simplicial sets.
\begin{definition}
    (See \cite{MR1650134}, Definition 4.2.18) Given a monoidal model category $\cC$, a $\cC$-\textit{model category}  is a $\cC$-module $\cD$(see Definition 4.1.6 for more details) with a model structure making $\cD$ into a model category such that the following conditions hold.
    \begin{enumerate}
        \item The action map $\bigotimes:\cC \times \cD \rightarrow \cD$ is a Quillen bifunctor.
        \item If $QS \xrightarrow{q} S$ is the cofibrant replacement for $S$ in $\cC$, then the map $X \bigotimes QS  \xrightarrow{1\bigotimes q} X \bigotimes S$ is a weak equivalence for all cofibrant $X$ $\in \cD$.
    \end{enumerate}
\end{definition}
\begin{definition}\label{def:pointedsimpmodcat}
    An $SSet$-model category is called a \textit{simplicial model category}. a pointed simplicial model category is the same thing as $SSet_{*}-$model category.
\end{definition}
We recall a Remark which concludes that the simplicial suspension and loop functor( i.e. the pointed tensor and cotensor with simplical circle $S^{1}=\triangle[1]_{+}/ \partial \triangle[1]_{+1})$ is a Quillen pair.
\begin{remark}
    (See \cite{MR1650134}, Remark 4.2.3) Suppose $\bigotimes:\cC \times \cD \rightarrow \mathcal{E}$ is a Quillen bifunctor. Then, if C is cofibrant, the functor $C \bigotimes -: \cD \rightarrow \mathcal{E}$ is a Quillen functor with right adjoint $Hom_{l}(C,-)$. Similarly if D is cofibrant , the functor $- \bigotimes D$ is a Quillen functor with right adjoint $Hom_{r}(D,-)$. Dually we have a statement with fibrant objects which we don't need.
\end{remark}
\begin{corollary}
    
Since a pointed simplical circle is a cofibrant object in $SSet_{*}$, so we get that the simplical suspension and loop functor is a Quillen pair on a pointed simplicial model category.
\end{corollary}
By the techniques of $framing$ in (See \cite{MR1650134}, Ch-5) Hovey shows that the homotopy category of a model category is naturally a closed $Ho SSet$-module and the homotopy category of a pointed model category is naturally a closed $Ho-SSet_{*}$-module.
And as noted in section 5.7 corollary 5.7.7 that a pointed simplicial model structure induces this closed $Ho-SSet_{*}$-module structure on the homotopy category of a pointed simplicial model category via the homotopy pseudo-2-functor.
\\This lets us give the following definition for homotopy category of any pointed model category. 
\begin{definition}
    (See \cite{MR1650134}) Definition 6.1.1) Suppose $\cC$ is a pointed model category. The \textit{suspension functor} $\Sigma: Ho \cC \rightarrow Ho\cC$ is the functor $X \mapsto X \wedge^{L} S^1$ defined by the closed action of $HoSSet_{*}$ on Ho$\cC$ as given in section 5.7. Dually the loop functor $\Omega: Ho \cC \rightarrow Ho\cC$ is the functor $X \mapsto RHom_{*}(S^{1},X)$.
\end{definition}
\begin{corollary} \label{cor:Quiequvforsimpstab}
    For a pointed simplicial model category the definition 2.16 is the one given by deriving the Quillen adjunction of corollary 2.15.
\end{corollary}

\subsection*{Pre-triangulated structure on the homotopy category of a model category}
We now recall the definition of cofiber and fiber sequences.

\begin{definition}
    (See \cite{MR1650134}, Definition 6.2.6)\label{def:cofseq} Suppose $\cC$ is a pointed model category. A \textit{cofiber sequence} in Ho $\cC$ is a diagram $X \xrightarrow{f} Y \xrightarrow{g} Z$ in Ho$\cC$ together with a right coaction of $\Sigma X$ on $Z$ which is isomorphic in Ho$\cC$ to a diagram of the form $A \rightarrow B \rightarrow C$ where $f$ is a cofibration of cofibrant objects in $\cC$ with cofiber $g$ and where $C$ has the right $\Sigma A$-coaction(This will give a coboundary map as in defintion 2.19 below) given by theorem 6.2.1. Dually a \textit{fiber sequence} is a diagram $X \rightarrow Y \rightarrow Z$ together with a right action of $\Omega Z$ on $X$ which is isomorphic to a diagram $F \xrightarrow{i} E \xrightarrow{p} B$ where $p$ is a fibration of fibrant objects with fiber $i$ and where $F$ has the right $\Omega B$-action given by theorem 6.2.1.
\end{definition}
\begin{definition} \label{def:boundarymap}
    (See \cite{MR1650134}, Definition 6.2.7) Suppose $\cC$ is a pointed model category, and $X \xrightarrow{f} Y \xrightarrow{g} Z$ in Ho$\cC$. The \textit{boundary map} is the map $\partial: Z \rightarrow \Sigma X$ in Ho$\cC$ which is the composite $Z \rightarrow Z\bigsqcup \Sigma X \xrightarrow{0 \times 1} \Sigma X$ \\ where the first map is the coaction as mentioned above. Dually, if $X \xrightarrow{f} Y \xrightarrow{g} Z$ is a fiber sequence, the \textit{boundary map} is the map $\partial :\Omega Z \rightarrow X$ which is the composite $\Omega Z \xrightarrow{(0,1)} X \times \Omega Z \xrightarrow{0 \times 1}  X$
\end{definition}
We now define what is called a pre-triangulated structure on a closed $HoSSet_{*}$-module, hence on the homotopy category of a pointed model category.
\begin{definition}
    (See \cite{MR1650134}, Definition 6.5.1) \label{def:preTristruct}Suppose $\mathcal{S}$ is a nontrival (right) closed $Ho SSet_{*}$-module. A \textit{pre-triangulation} on $\mathcal{S}$ is a collection of \textit{cofiber sequences}, or \textit{left triangles}, and \textit{fiber sequences}, or \textit{right triangles}, satisfying the following conditions.
    \begin{enumerate}
         
    
        \item A cofiber sequence is in particular a diagram of the form $X \xrightarrow{f} Y \xrightarrow{g} Z$ together with a right coaction of the cogroup $\Sigma X$ on $Z$. A fiber sequence is in particular a diagram of the form $X \xrightarrow{f} Y \xrightarrow{g} Z$ together with a right action of the group $\Omega Z$ on $X$.
        \item Every diagram isomorphic to a cofiber sequence is a cofiber sequence. Similarly, every diagram isomorphic to a fiber sequence is a fiber sequence. The isomorphisms respect the action as in \cite{MR1650134}, lemma 6.3.1.
        \item For any $X$, the diagram $* \rightarrow X \xrightarrow{1} X$ is a cofiber sequence(with the only possible coaction). The diagram $X \xrightarrow{1} X \rightarrow *$ is a fiber sequence (with the only possible coaction).
        \item Every map is a part of a cofiber fiber sequence, as in \cite{MR1650134}, Lemma 6.3.3.
        \item Cofiber sequences can be shifted to the right, and fiber sequences can be shifted to the left, as in \cite{MR1650134}, Proposition 6.3.4.
        \item Fill-in maps exist, as in \cite{MR1650134}, Proposition 6.3.5.
        \item Verdier's octahedral axiom and its dual both hold, as in \cite{MR1650134}, Proposition 6.3.6.
        \item Cofiber and fiber sequences are comptaible, as in \cite{MR1650134}, Proposition 6.3.7.
        \item The smash product preserves cofiber sequences in each variable. The functor $Hom(-,-)$ preserves fiber sequences in the second variable and converts cofiber sequences in the first variable into fiber sequences. The functor $Map(-,-)$ preserves fiber sequences in the second variable and converts cofiber sequences in the first variable into fiber sequences. See \cite{MR1650134}, corollary 6.4.2 and \cite{MR1650134}, Proposition 6.4.5 for exact statements.
    \end{enumerate}
\end{definition}
\begin{definition}
    (See, \cite{MR1650134}, Definition 6.5.2)\label{def:pretriCat} A \textit{pre-triangulated category} is a nontrivial closed $HoSSet_{*}$ modules $\mathcal{S}$ with all small coproducts and products, together with a pre-triangulation on $\mathcal{S}.$
\end{definition}
\begin{remark}
    As we note this \textit{pre-triangulated structure} exists on the homotopy category of a model category and a given Quillen adjunction say $(F \dashv G)$ induces an \textit{exact} adjunction $(LF \dashv RG)$between homotopy categories, where $LF$ preserves cofiber sequences and $RG$ preserves fiber sequences as in definition (\cite{MR1650134}, Definition 6.5.4). 
\end{remark}
 We now define a Stable model category.
 \begin{definition}\label{def:Hoveytricat}
     (See, \cite{MR1650134}, Definition 7.1.1) A \textit{triangulated category} is a pre-triangulated category in which the suspension functor $\Sigma$ is an equivalence of categories.\label{def:stabmodcat} A pointed model category is \textit{stable} if it's homotopy category is triangulated.
 \end{definition}
 \begin{example}
     All the categories in this paper such as $Ch(\mathcal{A})$, where $\mathcal{A}$ is any Grothendieck abelian category, \cite{MR1417719} $S$-modules, etc. are stable.
 \end{example}
 \begin{remark} \label{rem:factstri}
 \begin{enumerate}
     \item \label{ite:stabquilleneq} A simplicial model category is stable if it is pointed and the Quillen adjunction of corollary 2.15 is a Quillen equivalence.
     \item \label{ite:hoveyvsverdier} By proposition (\cite{MR1650134}, 7.1.6) a triangulated category in the sense of definition 2.23 with \textit{suspension functor} and \textit{cofiber sequences} becomes a classical triangulated category in the sense of \cite{neeman:trcat}.
     \item Although \cite{MR1417719} $S$-modules is also an example of simplicial stable cofibrantly generated model category. In folklore $Ch(\mathcal{A})$ can also be made simplicial. But we note that we don't need simplicial for our main result Theorem 3.5 and hence not for $Ch(\mathcal{A})$ as well.
\end{enumerate}
 \end{remark}



\subsection{Transferred model structure on module categories.}
%\todoinline{Cofibrantly generated transferred model structure}
We now state a result, which makes sure $\module{M}$ admits a \textit{right-induced} model structure from $\cC$.
Define a map of $\module{M}$ to be a \textit{weak equivalence} or \textit{fibration} if the underlying map inn $\cC$ is a \textit{weak equivalence} or \textit{fibration} respectively. Define a map in $M$-mod to be a \textit{cofibration} if it has the left lifting property with respect to all acyclic fibrations.
A \textit{regular $I$-cofibration} is a map in the class of the (possibly transfinite) compositions of pushouts of maps in $I$.

Let $X$ be a $M$-module. Define a \textit{path object} for $X$ to be an object in $M$-mod $X^{I}$ together with $M$-algebra maps $X \xrightarrow{\sim} X^{I} \rightarrow X \times X$ factoring the diagonal map, such that the first map is a weak equivalence and the second map is a fibration in the underlying category $\cC$.

The following is lemma 2.3 from \cite{MR1734325} (\emph{Caveat:} In that paper, $M$-algebras are what we call $M$-modules here).
\begin{lemma} \label{lem:schwship}
    Let $\mathcal{C}$ be a cofibrantly generated model category, with $I$ a set of generating cofibrations and $J$ a set of generating acyclic cofibrations. Let $(M,\mu,\eta)$ be a monad on $\mathcal{C}$ such that $M:\mathcal{C} \rightarrow \mathcal{C}$ commutes with filtered direct limits. Let $I_{M}$ and $J_{M}$ denote the images of $I$ and $J$ under the free $M$-module functor $F_{M}$. Assume that the domains and codomains of the morphisms in $I_{M}$ and $J_{M}$ are small with respect to regular $I_{M}$-cofibrations and regular $J_{M}$-cofibrations respectively. Suppose that  
    \begin{enumerate}
        \item every regular $J_{M}$-cofibration is a weak equivalence, or
        \item every object of $\cC$ is fibrant and every $M$-module admits a path object.
    \end{enumerate}
    Then the category M-mod is a cofibrantly generated model category with $I_{M}$ a generating  set of cofibrations and $J_{M}$ a generating set of acyclic cofibrations,
\end{lemma}
 The forgetful functor $G_{M}:M-mod \rightarrow \cC$ is said to \textit{right-induce} a cofibrantly generated model category structure on $M-mod$.
 \begin{remark} \label{rem:postcofib}
     \begin{enumerate}
         \item As mentioned in \cite{MR1734325} remark 2.4, when a category $\cC$ is a \textit{locally presentable category} every object of it is small. Further, when the category is not locally presentable then we can rely on the assumption that if underlying functor $M$ commutes with filtered direct limits, then so does $G_{M}$. Hence by adjointness, if every object of $\cC$ is small relative to $\cC$, then every free $M$-module is small relative to $M-mod$, so the smallness assumptions of lemma 2.26 holds. 
         \item \label{rem:simppath} As noted in \cite{MR1734325} remark 4.5, if $\cC$ is a simplical model category on which a simplical monad acts then M-mod obtains a simplicial structure in which case it always has path objects. In cases such as \cite{MR1417719} $S$-modules all objects are fibrant as well.
         \item Another set of conditions for which lemma 2.26 holds is presented in lemma \ref{lem:cndforsc} below.
     \end{enumerate}
 \end{remark}

\subsection{Algebraic Geometry}
\begin{notation*}
The Grothendieck abelian category in this paper $\mathcal{A}$ is the category of \textit{quasi-coherent sheaves} over a scheme $S$ denoted by $QCoh(S)$. We denote by $\Ch(QCoh(S))$ the unbounded chain complexes of quasi-coherent sheaves and it's homotopy category the derived category is denoted by $D(QCoh(S))$.

We now recall a model structure in Algebraic geometry called \textit{locally free model structure} as presented by Hovey \cite{MR1814077} on $\Ch(QCoh(S))$ for a suitable scheme $S$ as below. \\ 
Let $\mathcal{A}$ be a Grothendieck category and recall that every object in a Grothendieck category is small. This lets us deal with smalless in transferred model structures as in Remark 2.27 (1) above.

Recall that $S^{n-1},D^{n}:\mathcal{A} \rightarrow Ch(\mathcal{A})$ are the functors that send $A \mapsto S^{n-1}A$ i.e. the complex whose only non-zero object is $A$ in degree $n-1$ and $A \mapsto D^{n}A$ the complex whose only nonzero objects are $A$ in degree $n-1$ and $n$ respectively.
\end{notation*}

Recall the following definition from \cite{MR1814077}.
\begin{definition}
    (\cite{MR1814077}, Definition 1.1) \label{def: I and J}Suppose $\mathcal{M} $ is a set of monomorphisms in a Grothendieck category $\mathcal{A}$. Let $\mathcal{F}$ denote the set of codomains of the maps of $\mathcal{M}.$ We will say that $\mathcal{M}$ is pointed if $0 \in \mathcal{F}$ and, if $F \in \mathcal{F}$, then $0 \rightarrow F$ is in $\mathcal{M}.$ Define $J$ to be the set of all $D^{n}f$, where $n$ is an integer and $f \in \mathcal{M}$. Define $I$ to be the union of $J$ and the maps $S^{n-1}F \rightarrow D^{n}F$ for $F \in \mathcal{F}$ and $n$ an integer. Then define a map $p$ to be an $\mathcal{M}\text{-fibration}$ if $p$ is in $J\text{-inj}$, and define $p$ to be an $\mathcal{M}$-cofibration if $p$ is in $I\text{-cof}$. 
\end{definition}

Recall that an object $B$ of a Grothendieck category $\mathcal{A}$ is said to have \textit{finite projective dimension} if there is an integer $n_{0}$ such that $Ext_{\mathcal{A}}(B,C)=0$ for all $n \ge n_{0}$ and all objects $C$ of $\mathcal{A}$. 

We now recall a result about existence of a cofibrantly generated model structure on $Ch(\mathcal{A})$ .
\begin{theorem} (\cite{MR1814077} Theorem 2.2)
    Suppose $\mathcal{A}$ is a Grothendieck category with a set $\mathcal{F}$ of generators, each of which has finite projective dimension. Let $\mathcal{M}$ denote the set of inclusions $A \rightarrow F$ of subobjects of objects $F \in \mathcal{F}$. Then there is a proper cofibrantly generated model structure on $Ch(\mathcal{A})$, where the weak equivalences are the quasi-isomorphims, the fibrations are the dimensionwise split surjections with dimensionwise injective kernel, and the cofibrations are the $\mathcal{M}$-cofibrations.
\end{theorem}

Also recall that a Noetherian scheme $S$ is said to have \textit{enough locally frees} if every coherent sheaf on $S$ is a quotient of a locally free sheaf of finite rank. %\todo{For example see Hartshorne Ex 3.6.8}

We now recall  \cite[proposition 2.3]{MR1814077} which allows the existence of the mentioned \textit{locally free model structure}.

\begin{proposition}[\cite{MR1814077}, Proposition 2.3] Suppose S is a noetherian scheme with enough locally frees. In addition, suppose that S is either finite-dimensional or is seperated. Then the locally free sheaves of finite rank are generators of finite projective dimension for the category QCoh(S) of quasi-coherent sheaves on S.
\end{proposition}

We have the following existence theorem as a corollary to Thereom 2.29 and proposition 2.30, which is \cite{MR1814077} Theorem 2.4.
\begin{theorem} (\cite{MR1814077} Theorem 2.4)\label{thm:locallyfreeModelStruct}
    Suppose $S$ is a noetherian scheme with enough locally frees, and suppose that $S$ is either finite-dimensional or seperated. Then there is a proper,cofibrantly generated, model structure on the category $\Ch(QCoh(S))$ of unbounded complexes of quasi-coherent sheaves, where the weak equivalences are the quasi-isomorphisms and the fibrations are the dimensionwise split surjections with dimensionwise injective kernel.
\end{theorem}

\begin{remark}
\begin{enumerate}
    \item If $f:X \rightarrow Y$ is a map  between schemes satisfying the hypotheses of thereom 2.31 then $f^{*}: QCoh(Y) \rightarrow QCoh(X)$ will induce a Quillen functor between the locally free model structures if and only if $f^{*}$ is exact. Note that for converse $f^{*}$ preserves locally free sheaves of finite rank.
    
    \item We have the following two deductions about $D(QCoh(S))$ from the locally free model structure: 
    \begin{enumerate}
        \item \cite[corollary 2.5]{MR1814077} Suppose $S$ is a noetherian scheme with enough locally frees, and $S$ is either finite-dimensional or seperated. Then the locally free sheaves of finite rank form a set of small weak generators for the $D(QCoh(S))$.
        \item \cite[corollary 2.6]{MR1814077} Suppose $S$ is a finite-dimensional noetherian scheme with enough locally frees. Then the category $D(QCoh(S))$ is a unital algebraic stable homotopy category, where the generators are the locally free sheaves of finite rank.             
    \end{enumerate}  
\end{enumerate}

\end{remark}






\section{Stable Compatibility} \label{sec:stablecpt}
In this paper, we assume the following:
\begin{assumption} \label{asm:stabcomp}
    \begin{enumerate}
        \item The category $\cC$ is pointed.
        \item We work with monads $(M, \mu, \eta)$ (see \ref{ssc:monads}), where the underlying endofunctor $M: \cC \rightarrow \cC$ of the monad preserves the zero object, and hence, the category of $M$-modules $\module[\cC]{M}$ is also pointed. The free functor $F_M$ and its right adjoint $G_M$ map the zero object in the appropriate category to the zero object
        of the other.
        \item \label{itm:suslooqe} Moreover, $\cC$ is a cofibrantly generated stable model category such that $(\Sigma \dashv \Omega)$ is a Quillen equivalence. This is the case, for example, when the cofibrantly generated stable model category is simplicial as noted in remark \ref{rem:factstri}\eqref{ite:stabquilleneq}.
        \item \label{itm:suscomGM} We also assume that $M$ commutes with the suspension $\Sigma$.
        \item \label{itm:rtmodstr} Moreover, when $\module{M}$ admits a right-transferred model structure, we shall assume that 
        $\adjunction{F_M}{\cC}{\module[\cC]{M}}{G_M}$ is a Quillen adjunction. The suspension and loop functors
        in $\cC$ induces a suspension functor $\Sigma_{\module{M}}$ and a loop functor $\Omega_{\module{M}}$ on $\module{M}$.
        We shall assume that, these induced suspension and loop functors commute with $G_M$.
    \end{enumerate}
\end{assumption} 

\begin{lemma}
    Under the assumptions \ref{asm:stabcomp}, 
    \[
      \adjunction{\Sigma_{M-mod}}{\module{M}}{\module{M}}{\Omega_{M-mod}}
    \]
    is a Quillen pair.
\end{lemma}
 \begin{proof}
     That $\Omega_{\module{M}}$ is a right Quillen functor is checked easily as follows. 
     We need to show that $\Omega_{\module{M}}$ preserves fibrations and acyclic fibrations.
     Let $f$ be a fibration (respectively, an acyclic fibration) in $\module{M}$. 
     Since both fibrations and weak equivalences in $\module{M}$ are created via $G_M$, 
     we need to check that $G_{M}(\Omega_{\module{M}}f)$ is a fibration (respectively, an acyclic fibration)
     in $\cC$. Since  $\Omega$ is a right Quillen functor on $\cC$ and since $G_{M}\Omega_{\module{M}}=\Omega G_{M}$, 
     $\Omega_{\module{M}} $ is a right Quillen functor. 
     
     Thus $(\Sigma_{\module{M}} \dashv \Omega_{\module{M}})$ forms a Qullen adjuntion (see, eg,
     \cite[lemma 1.3.4]{MR1650134}).
 \end{proof}

The following lemma describes a criterion for the Quillen adjuntion $(\Sigma_{\module{M}} \dashv \Omega_{\module{M}})$ to be a
Quillen equivalence.

\begin{lemma} \label{lem:shloopqe}
    Continuing with the same setup as in the previous lemma, if  $G_{M}$  preserves cofibrant objects 
    then the Quillen adjunction $\Sigma_{\module{M}} \dashv \Omega_{\module{M}}$ is a Quillen equivalence. 
\end{lemma}
\begin{proof}
  Note that $\Sigma_{\module{M}}$ is a Quillen equivalence \cite[definition 1.3.12]{MR1650134} if a map
  $\Sigma_{\module{M}}X \rightarrow Y$ is a weak equivalence in $\module{M}$ if and only if the corresponding map 
  $X \rightarrow \Omega_{\module{M}}Y$ is a weak equivalence in $\module{M}$ for all cofibrant objects $X$ and
  fibrant objects $Y$ in $\module{M}$.

  Recall that weak equivalences in $\module{M}$ with the transferred model structure are created by weak equivalences in $\cC$
  via $G_{M}$ and that $G_{M}\Sigma_{\module{M}}=\Sigma G_{M}$.  
  Let $x = G_M X$ and $y = G_M Y$. Since $G_{M}$ is a right Quillen functor, $y$ is a fibrant object in $\cC$,
  and by our hypothesis, $x$ is a cofibrant object in $\cC$.
  Therefore, we are reduced to showing that $\Sigma x \rightarrow y$ is a weak equivalence in $\cC$ if and only if
  the corresponding map $x \rightarrow \Omega y$ is a weak equivalence in $\cC$. This already holds as 
  $\Sigma \dashv \Omega$ is a Quillen equivalence on $\cC$.
\end{proof}

\begin{definition} \label{def:stabcomp}
  A monad $M$ satisfying the assumptions \ref{asm:stabcomp} is called \emph{stable compatible}
  (or, \emph{$M$ satisfies stable compatibility}) if the corresponding forgetful functor $G_M$
  preserves cofibrant objects.    
\end{definition}

We now highlight the importance of stable compatibility in the following theorem.
\begin{theorem} \label{the:stabcomp}
  Let $\cC$ be a stable cofibrantly generated model category.
  Suppose $M$ is a monad on $\cC$ which satisfies the following.
  \begin{enumerate}
    \item
      The category $\module{M}$ admits a right transferred model structure, which is 
      cofibrantly generated, and such that the induced adjunction $\Sigma_{\module{M}}
      \dashv \Omega_{\module{M}}$ is a Quillen adjunction. 
    \item The monad $M$ satisfies stable compatibility.
  \end{enumerate}
  Then $\module{M}$ is a stable model category. Moreover, if the monad is induced by a
  Quillen adjunction $\adjunction{F}{\cC}{\cD}{G}$, the underlying comparison functor $K$
  is such that $K \circ F=F_{M}$ and $G=G_{M} \circ K$ as in the diagram below   
  \[
    \begin{tikzcd}
      & \cC  \arrow[dr,"F_{M}", shift left=1mm] \arrow[d,"F", shift left=1mm] \\
      & \cD \arrow[dashrightarrow,r,"\exists !K" below] \arrow[u,"G",shift left=1mm]& \module{M} \arrow[ul, "G_{M}",shift left=1mm] 
    \end{tikzcd}
  \]
  and it preserves fibrations and acyclic fibrations.
\end{theorem}

\begin{proof}
  That $\module{M}$ is a stable model category is immediate from our previous discussion (see lemma \ref{lem:shloopqe}). 
   
   The existence of $K$ follows as usual. To see that $K$ preserves fibrations (respectively, acyclic fibrations),
   we use the fact that $G_{M}$ creates fibrations and weak equivalences in a right transferred model structure.
   Let $f$ be a fibration (respectively, acyclic fibration) in $\module{M}$. Since $G$ is a right Quillen functor, 
   $Gf$ is again a fibration (respectively acyclic fibration). Now, $G=G_{M} \circ K$ implies that $Kf$ is a 
   fibration (respectively, acyclic fibration).  
\end{proof}

\begin{remark} \label{rmk:stabcomp}
\begin{enumerate}
    \item The categorical diagram involving the Kleisli category as in Proposition \ref{pro:kleisliem} does not extend in this scenario, 
      as it is not a complete-cocomplete category.
    \item It is not true in general that $G_{M}$ commutes with suspension, i.e. $G_{M}\Sigma_{M-mod} \neq \Sigma G_{M}$ (however, see remark \ref{rem:gmstable} below).
    \item \label{itm:simmonad} It is a well-known result that if $\cC$ is a simplicial,  cofibrantly generated model category 
    with an action of a simplicial monad, which satisfies the assumptions of Schwede and Shipley \cite[lemma 2.3]{MR1734325},
    then $\module{M}$ is a simplicial cofibrantly generated model category with 
    $(F_{M} \dashv G_{M})$ being a simplicial Quillen adjunction (see Remark \ref{rem:postcofib}\eqref{rem:simppath}). The condition of stable compatibility
    extends this to a stable model structure.
    \item  \label{itm:chRasamc} 
    The converse to above theorem is not true in general; that is, the module category
    being stable does not imply that the monad is stable compatible. We give the
    following example.
    
    Consider the case of a commutative ring $R$ and an $R$-algebra $S$. We consider the projective/standard model structure on $Ch(R)$, the category of unbounded chain complexes of $R$-modules, as described in \cite[2.3.11]{MR1650134}. It is known to satisfy the monoid axiom of Schwede and Shipley (see \cite[section 5, \textit{Chain complexes}]{MR1734325}).
    Consider the monoid object given by the discrete DGA $S_{\bullet}= \cdots 0 \rightarrow S \rightarrow 0 \cdots$ in $Ch(R)$, where $S$ is situated at degree $0$.
    By definition, it follows that the Eilenberg-Moore category of 
    $S_{\bullet}$-modules in $Ch(R)$ is isomorphic to the category $Ch(S)$;
    the transferred model structure on the monad $T=(S_{\bullet} \bigotimes \_)$ 
    is the same as projective model structure on $Ch(S)$;
    and the forgetful functor $G_{M}$ is just degreewise restriction of scalars.
    Moreover, the suspension functor is the shift functor in this case
    and it commutes with forgetful functor. Further, it is a Quillen equivalence on both categories. By \cite[2.3.6]{MR1650134}, if $A_{\bullet}$ is a cofibrant object,
    each $A_{n}$ is a projective S-module for each $n$. However, restricting
    a projective $S$-module is not a projective $R$-module unless 
    $S$ is also a projective $R$-module. Thus $G_M$ does \emph{not} preserve cofibrant
    objects and stable compatibility fails.
\end{enumerate}
\end{remark}

\subsection{Some criteria for stable compatibility}
We would now like to point out some criteria under which the hypothesis of theorem \ref{the:stabcomp}. We'll use this to show that stable compatibility holds in the applications stated in this paper. The idea is to allow $G_{M}$ to have a good homotopical control. We state it in the form of the following lemma.
\begin{lemma} \label{lem:cndforsc}
    If the underlying endofunctor of the monad $M:\cC \rightarrow \cC$ preserves colimits, generating cofibrations, and generating acyclic cofibrations, then the monad satisfies hypothesis of theorem \ref{the:stabcomp}.
\end{lemma}
\begin{proof}
Let $I$ be the set of generating cofibrations, and $J$ the set of generating acyclic cofibrations of $\cC$. The candidate for
respective generating sets in a transferred model structure are exactly $F_{M} I$ and $F_{M} J$. Now, to transfer the right-induced
model structure, we check the conditions of Schwede and Shipley, as in lemma \ref{lem:schwship}.

Let $f: X \rightarrow Y$ be an $F_{M}(J)$-cell complex as in definition \ref{def:I-cell}, we need to show that $f$ is a weak equivalence. This is same as showing that
$G_{M} f$ is a weak equivalence. Let $f$ be the $\lambda$-transfinite composition of the following pushouts
\begin{equation} \label{eqn:pushoutf}
    \begin{tikzcd}
       \bigsqcup F_{M}A \arrow{r} \arrow[swap]{d} & X_{\beta} \arrow{d} \\
        \bigsqcup F_{M}B \arrow{r} & X_{\beta +1}
     \end{tikzcd}
\end{equation} 
where $X_{0}=X$.

Recall that $M=G_{M}F_{M}$ and $M$ preserves colimits (by assumption). Also $G_M$ creates, and hence preserves, colimits (cf \cite[theorem 5.6.12]{MR4727501} and \cite{MR2178101}).
Applying $G_{M}$ to the diagram \eqref{eqn:pushoutf}, we get the following $M(J)$-cell complexes in $\cC$
\begin{equation} \label{eqn:MJccmplx}
    \begin{tikzcd}
       \bigsqcup MA \arrow{r} \arrow[swap]{d} & G_{M}X_{\beta} \arrow{d} \\
        \bigsqcup MB \arrow{r} & G_{M}X_{\beta +1}
\end{tikzcd}
\end{equation}
Since $M$ preserves $J$, diagram \eqref{eqn:MJccmplx} is in fact a $J$-cell in $\cC$
%\todo{Assure that a $\lambda$-sequnece stays a valid sequnece set-theoretically}
and $G_{M} f$ is a $J$-cell in $\cC$ hence an acyclic cofibration and hence a weak equivalence.

Similarly, we can start with $X_{0}=*$ in $F_{M}I$-cell complexes and since $M$ preserves $I$ and $G_{M}$ preserves retract, we get the stable
compatibility of $M$.
\end{proof}

\begin{corollary} \label{cor:prescofb}
    The forgetful functor $G_{M}$ preserves cofibrations.
\end{corollary}
\begin{proof}
    Let $f:X \rightarrow Y$ be a cofibration in $\module{M}$. The proof of lemma \ref{lem:cndforsc} shows that $G_{M}(F_{M}I\text{-cell complex}) \subset I\text{-cell complex}$ in $\cC$. Since a cofibration in $\module{M}$ is a retract of an $F_{M}I$-cell complex and $G_{M}$ preserves both retracts and
    cell complexes, $G_{M}f$ is a retract of an $I$-cell complex in $\cC$. Hence, $G_M f$ is a cofibration.
\end{proof}
\begin{remark} \label{rem:gmstable}
    It is straightforward to check assumption \ref{asm:stabcomp}\eqref{itm:suscomGM} under hypothesis of \ref{lem:cndforsc}, 
    provided $\Sigma_{\module{M}}$ is well-defined and makes a Quillen pair with $(\Sigma_{\module{M}} \dashv \Omega_{\module{M}})$. 
    By definition, $\Sigma_{\module{M}}$X is the homotopy pushout of the diagram
    \[
    \begin{tikzcd}
      X \arrow{r} \arrow[swap]{d} & *  \\
       * 
    \end{tikzcd}
    \] 
   which is the ordinary pushout but with the above diagram replaced by a cofibrant one. 
   By corollary \ref{cor:prescofb}, $G_{M}$ preserves cofibrations, and we know that it preserves push out diagrams. Thus, it follows that
   $G_M$ commutes with suspension.
\end{remark}





\section{Separable monads} \label{sec:sepmonad}

In \cite{balmer:septrcat}, Balmer showed that for a separable monad on a tensor triangulated category, the corresponding
category of modules $\module{M}$ admitted a triangulated structure, where the corresponding functor $G_M$ generated the
distinguished triangles. In this section, we study the effect of introducing a condition like separability in our model category set up.

We would like to recall the difference between the notion of pretriangulated in Hovey \cite{MR1650134} and Balmer \cite{balmer:septrcat}
to avoid confusion. For Hovey, a pretriangulated structure on a homotopy category of a model category $\cC$ means a collection of 
cofiber sequences and fiber sequences, which satisfy the
axioms as in definitions \ref{def:preTristruct} and \ref{def:pretriCat}. Such a structure always exists on the homotopy category of a 
cofibrantly generated model category (see remark \ref{def:Hoveytricat}). He then defines a triangulated category as a pretriangulated category when 
the suspension $\Sigma$ is an equivalence of homotopy category as in definition \ref{def:Hoveytricat}. Hovey observes that these triangulated categories are triangulated
in the classical sense described in the next paragraph.

For Balmer, a pretriangulated category means collection of distinguished triangles in a category satisfying all the axioms of
a triangulated category, except the octahedral axiom. See for example, \cite[Chapter 1, section 1.1]{neeman:trcat}. Here, the suspension functor
is always an autoequivalence (see definition \ref{def:monad}). A triangulated category \cite[definition 1.3.13 and proposition 1.4.6]{neeman:trcat} is a pretriangulated category where the octahedral axiom \cite[remark 1.4.7]{neeman:trcat} holds.

In this section, we study Balmer's treatment of separability for triangulated categories coming from the model categorical framework.
As mentioned in the introduction, this provides more information on \cite[remark 6.1]{balmer:septrcat}. In our situation, the distinguished triangles are the cofiber sequences of 
definitions \ref{def:cofseq} and \ref{def:boundarymap} as explained in remark \ref{rem:factstri}\eqref{ite:hoveyvsverdier}.
        
In the case of triangulated categories \cite{balmer:septrcat}, suspension functor on $\cC$ induce a suspension functor 
$\module[\cC]{M}$, which turns out to be autoequivalent. Contrast with this the situation involving a suitable model category $\cC$ with a monad $M$
satisfying stable compaibility. In this case, the induced suspension functor on $\Ho(\module[\cC]{M})$ is not automatically
an equivalence. It happens to be an autoequivalence, for example, when the suspension functor on the model category $\module[\cC]{M}$
is a Quillen equivalence. Hence, stable compatibility is much weaker than the assumptions in \cite{balmer:septrcat}.

Balmer further works with idempotent complete categories. We continue with the assumption that the categories
involved are idempotent complete in order to recover the work. 

\begin{definition}
    A model category is \emph{homotopy idempotent complete} if its homotopy category is an idempotent-complete category.
\end{definition}


\begin{definition}
    A monad $(M, \eta,\mu)$ is called \textit{separable} if the multiplication $\mu \colon M \circ M \rightarrow M$ admits section 
    $\sigma \colon M \rightarrow M \circ M$, that is $\mu \circ \sigma = \id_M$, and such that $M \mu \circ \sigma M = \sigma \circ \mu
    = \mu M \circ M \sigma$.
\end{definition}
\begin{remark}
  Note that, by Balmer \cite[section 3, \S\S 3.5 -- 3.11]{balmer:septrcat}, $M$ is separable, if and only if $G_M$ is a separable functor, which is equivalent to saying that the counit $\epsilon \colon F_{M}G_{M} \rightarrow \id_{\module{M}}$ admits a section $\upxi: \id_{\module{M}} \rightarrow F_{M}G_{M}$.
\end{remark}

\begin{assumption} \label{asm:sepassmpt}
For the rest of this section, we assume that the separable  monad also satisfies the assumptions of theorem \ref{the:stabcomp},
for example, via lemma \ref{lem:cndforsc} or remark \ref{rmk:stabcomp}\eqref{itm:simmonad}.
\end{assumption}

\begin{example} \label{exa:sepalgeb}
    An example of such a monad comes alike from remark 
    \ref{rmk:stabcomp}\eqref{itm:chRasamc} when the $R$-algebra $S$ is finite \'etale . The forgetful functor $G_{M}$ becomes a 
    bi-Quillen functor and, in particular, satisfies stable compatibility. As an aside, also observe that in this
    case the conditions of lemma \ref{lem:cndforsc} are satisfied, giving us another way to check stable compatibility.
    The sections apply degreewise to get a 
    separable monad $S \bigotimes \_$ on $\Ch(R)$ with the standard/projective model structure mentioned in 
    \ref{rmk:stabcomp}\eqref{itm:chRasamc}.
\end{example}

We begin with a lemma that reflects the interaction of a separable monad with the transferred model structure.
We work with maps between bifibrant objects in $\module{M}$. For a map $f$ in the model category,, the corresponding map in the homotopy category
is denoted by $[f]$.  Using the small object arguement, factorize $f$ as $X \xrightarrow{i_{f}} E_{f} \xrightarrow{p_{f}}Y$, where $i_{f} \in F_{M}I\text{-cell}$ and $p_{f} \in RLP(F_{M}I)$, where $RLP(F_M I)$ denotes all the maps satisfying right lifing property with respect to maps in $F_M I$. 
As $X$ is cofibrant, the object $E_{f}$ is cofibrant  and as $Y$ is fibrant, $E_{f}$ is fibrant too. Therefore $E_{f}$ is bifibrant.

\begin{lemma} \label{lem:iftoFiGf}
    Under assumption \ref{asm:sepassmpt}, if $M$ is separable, then there exists a map in the arrow 
    category $i_{f} \rightarrow F_{M}i_{G_{M}f}'$ .    
\end{lemma}
\begin{proof}
 Let us consider the factorisation of $G_{M}f$ in $\cC$. Applying the small object argument in $\cC$, we get $G_{M}X \xrightarrow{i'_{G_{M}f}}E'_{G_{M}f} \xrightarrow{p'_{G_{M}f}}G_{M}Y$, along with the following commutative diagram in $\cC$:
\begin{equation} \label{eqn:cofibar1}
    \begin{tikzcd}
       G_{M}X \arrow{r}{i'_{G_{M}f}} \arrow{r} \arrow[swap]{d}{G_{M}i_{f}} & E'_{G_{M}f} \arrow{d}{p'_{G_{M}f}} \\
        G_{M}E_{f} \arrow{r}{G_{M}p_{f}} & G_{M}Y.
     \end{tikzcd}
\end{equation} 
By corollary \ref{cor:prescofb}, the map $G_{M} i_{f}$ is a cofibration, and since the map $p'_{G_{M} f}$ is an acyclic fibration, by the RLP of acyclic fibrations with respect to cofibrations, we get a lift in the diagram \ref{eqn:cofibar1} which we denote by $t_{f}$:  
\begin{equation} \label{eqn:cofibar2}
    \begin{tikzcd}
       G_{M}X \arrow{r}{i'_{G_{M}f}} \arrow{r} \arrow[swap]{d}{G_{M}i_{f}} & E'_{G_{M}f} \arrow{d}{p'_{G_{M}f}} \\
        G_{M}E_{f} \arrow{r}{G_{M}p_{f}} \arrow[ur, dashed,"t_{f}"] & G_{M}Y
     \end{tikzcd}
\end{equation} 
 
Since both the maps $p'_{G_{M}f}$ and $G_{M}p_{f}$ are weak equivalences, by the 2-out-of-3 property of weak equivalences, so is $t_{f}$. Further, since $G_{M}X$ is a cofibrant object, so is $E'_{G_{M}f}$, therefore $t_{f}:G_{M}E_{f} \rightarrow E'_{G_{M}f}$ is a weak equivalence between cofibrant objects.

Since $F_{M}$ is a left Quillen functor, it preserves weak equivalences between cofibrant objects. Applying $F_{M}$ to the commutative diagram \eqref{eqn:cofibar2}, we get the following commutative diagram in $\module{M}$:
\begin{equation} \label{eqn:cofibar3}
    \begin{tikzcd}
       F_{M}G_{M}X \arrow{r}{F_{M}i'_{G_{M}f}} \arrow{r} \arrow[swap]{d}{F_{M}G_{M}i_{f}} & F_{M}E'_{G_{M}f} \arrow{d}{F_{M}p'_{G_{M}f}} \\
        F_{M}G_{M}E_{f} \arrow{r}{F_{M}G_{M}p_{f}} \arrow[ur, dashed,"F_{M}t_{f}"] & F_{M}G_{M}Y
     \end{tikzcd}
\end{equation} 

Now, by the upper-left triangle of diagram \eqref{eqn:cofibar3} and the naturality of sections
$\upxi_{X} \colon X \rightarrow F_{M}G_{M}X$ and $\upxi_{E_{f}} \colon E_{f} \rightarrow F_{M}G_{M}E_{f}$, 
the map $(\upxi_{X},(F_{M}t_{f}\circ \upxi_{E_{f}}))$ is the map claimed in the statement of the lemma:
\begin{equation}
    \begin{tikzcd}
       X \arrow[rr, "\upxi_{X}"]  \arrow[d, "i_{f}"'] & & F_{M}G_{M}X \arrow[d, "F_{M}i'_{G_{M}f}"] \\
       E_{f} \arrow[r, "\upxi_{E_{f}}"] 
       \arrow[rr, bend right, "F_M t_f \circ \upxi_{E_f}"'] & F_M G_M E_f \arrow[r, "F_{M}t_{f}"] & F_{M} E'_{G_{M}f}
     \end{tikzcd}
\end{equation} 
\end{proof}

\begin{remark}
We note that $X$, $E_{f}$, $F_{M}G_{M}X$, $F_{M}E'_{G_{M}f}$ are all cofibrant objects in $\module{M}$
and $i_{f}$ and $F_{M}i'_{G_{M}f}$ are cofibrations.
\end{remark}

Next, we want to extend the above map into a map of cofiber sequences. Before that, we recall a result from Hovey.
\begin{lemma}[Hovey \cite{MR1650134}, Proposition 6.2.5] \label{lem:extcofib}
Suppose $\cC$ is a pointed model category and we have a commutative square of cofibrant objects
\[
  \begin{tikzcd}
        A' \arrow{r}{f'} \arrow{d}{q_{1}} \arrow[swap]{d} & B' \arrow{d}{q_{2}} \\
        A \arrow{r}{f} & B
  \end{tikzcd}
\] 
where $f'$ and $f$ are cofibrations, with cofibers $g':B' \rightarrow C'$ and $g: B \rightarrow C$ respectively. 
Then the induced map $q_{3}:C' \rightarrow C$ is equivariant in $Ho(\cC)$ with respect to the cogroup homomorphism $\Sigma q_{1}$.

The corresponding statement for fibrations hold by duality.   
\end{lemma}
\begin{corollary} \label{cor:indcisom}
    If $q_{1}$ and $q_{2}$ are weak equivalences in lemma \ref{lem:extcofib} then the induced map between the
    cofiber sequences is an isomorphism.
\end{corollary}
\begin{proof}
Since the cofiber of  $f': A' \rightarrow B'$ map is the pushout of the diagram: 
\begin{equation} \label{eqn:diaapcof}
    \begin{tikzcd}
       A' \arrow{r}{f'}  \arrow[swap]{d} & B'  \\
        * 
     \end{tikzcd}
\end{equation} 

and similary the cofiber of $f: A \rightarrow B$:
\begin{equation} \label{eqn:diagacof}
    \begin{tikzcd}
       A \arrow{r}{f}  \arrow[swap]{d} & B \\
        * 
     \end{tikzcd}
\end{equation} 
    Then $(\id_{*},q_{1},q_{2})$ are the object-wise weak equivalences between the diagrams \eqref{eqn:diaapcof} and 
    \eqref{eqn:diagacof}. By the cube lemma \cite[lemma 5.2.6]{MR1650134}, the induced map $q_{3}: C' \rightarrow C$
    is a weak equivalence as well.
\end{proof}

\begin{lemma} \label{lem:dirincof}
The map $i_{f} \rightarrow F_{M}i_{G_{M}f}'$ in lemma \ref{lem:iftoFiGf} extends to an inclusion map
of the cofiber sequence $C_1$ associated with $i_{f}$ into the cofiber sequence $C_2$ associated with
$F_{M}i_{G_{M}f}'$ in $\Ho(\module{M})$. Moreover, the sequence $C_1$ is a direct summand of $C_2$ under this inclusion map.
\end{lemma}
\begin{remark}
    Before we proceed to give a proof, note that this lemma is not saying that cofiber sequence associated with $i_{f}$ in $\Ho(\module{M})$is a direct summand of $L F_{M}$ applied to the cofiber sequence of $i'G_{M}f$ in $\Ho \cC$. Nor is it implying that the triangle associated to $[f]$ is a direct summand of the $LF_{M}$ applied to the triangle associated to the map $RG_{M}([f])$. For the latter follow the proof of Balmer.
\end{remark}
\begin{proof}
We note that the map $i_{f} \rightarrow F_{M}i_{G_{M}f}'$ as defined factors as 
    \[\begin{tikzcd}
         X \arrow{r}{\upxi_{X}} \arrow[swap]{d}{i_{f}} & F_{M}G_{M}X \arrow{d}{F_{M}G_{M}i_{f}} \arrow{r}{id} & F_{M}G_{M}X \arrow{d}{F_{M}i'_{G_{M}f}}  \\
         E_{f} \arrow{r}{\upxi_{E_{f}}} & F_{M}G_{M}E_{f} \arrow{r} {(i_{2}=F_{M}t_{f})} & F_{M}E'_{G_{M}f} 
    \end{tikzcd}\]

By naturality of sections and counit, the first square \[\begin{tikzcd}
         X \arrow{r}{\upxi_{X}} \arrow[swap]{d}{i_{f}} & F_{M}G_{M}X \arrow{d}{F_{M}G_{M}i_{f}} \arrow[l,"\epsilon_{X}",shift left=1ex]   \\
         E_{f} \arrow{r}{\upxi_{E_{f}}} & F_{M}G_{M}E_{f}  \arrow[l,"\epsilon_{E_{f}}",shift left=1ex]
    \end{tikzcd}\]   

gives an idempotent of the square
\[\begin{tikzcd}
         F_{M}G_{M}X \arrow{r}{\upxi_{X}\epsilon_{X}} \arrow[swap]{d}{F_{M}G_{M}i_{f}} & F_{M}G_{M}X \arrow{d}{F_{M}G_{M}i_{f}}    \\
         F_{M}G_{M}E_{f} \arrow{r}{\upxi_{E_{f}}\epsilon_{E_{f}}} & F_{M}G_{M}E_{f}  
    \end{tikzcd}\] 
by lemma \ref{lem:extcofib}, this extends to an idempotent map of the cofiber sequence associated to the map $F_{M}G_{M}i_{f}$ and as $\Ho(\module{M})$ is an 
idempotent complete category, the cofiber sequence associated to $i_{f}$ is a direct summand of the former. Moreover, the pair 
$(\id_{F_{M}G_{M}X}, F_{M}t_{f})$ is a pair of weak equivalences, and by corollary \ref{cor:indcisom}, the cofiber sequence associated to $F_{M}G_{M}i_{f}$ is 
isomorphic to the cofiber sequence associated with the map $F_{M}i'_{G_{M}f}$ . So via the isomorphism the cofiber associated to $i_{f}$ becomes a direct summand of 
the cofiber associated to $F_{M}i'_{G_{M}f}$.
\end{proof}

\begin{remark} \label{rem:sepcompu}
Let us consider a pair of composable maps  $X \xrightarrow{f} Y \xrightarrow{g} Z$. Applying the functorial factorization, we obtain 
$i_{f}$, $i_{g}$ and $i_{gf}$ and since the factorization is functorial on the arrow category, the map $(\id_{X},g)$  
\[\begin{tikzcd}
        X \arrow{r}{f} \arrow{d}{\id_{X}}  & Y \arrow{d}{g} \\
        X \arrow{r}{gf} & Z
     \end{tikzcd}
     \]   
induces a map $s_{f}:E_{f} \rightarrow E_{gf}$ 
\[\begin{tikzcd}
        X \arrow{r}{i_{f}} \arrow{d}{\id_X}  & E_{f} \arrow{d}{s_{f}} \arrow{r}{p_{f}} & Y \arrow{d}{g}\\
        X \arrow{r}{i_{gf}} & E_{gf} \arrow{r}{p_{gf}} & Z
     \end{tikzcd}
     \]   
such that $s_{f}i_{f}=i_{gf}$ and $gp_{f}=p_{gf}s_{f}$ and similarly the map $(f,\id_{Z})$  
\[\begin{tikzcd}
        X \arrow{r}{gf} \arrow{d}{f}  & Z \arrow{d}{id_{Z}} \\
        Y \arrow{r}{g} & Z
     \end{tikzcd}
     \]  
     induces a map $s_{gf}:E_{gf} \rightarrow E_{g}$ giving a commutative diagram  
     \[\begin{tikzcd}
        X \arrow{r}{i_{gf}} \arrow{d}{f}  & E_{gf} \arrow{d}{s_{gf}} \arrow{r}{p_{gf}} & Z \arrow{d}{\id_{Z}}\\
        Y \arrow{r}{i_{g}} & E_{g} \arrow{r}{p_{g}} & Z.
     \end{tikzcd}
     \]
     Note that we obtain similar commutative diagrams in $\cC$ when we apply the functorial factorization to $G_{M}f$, 
     $G_{M}g$ and $G_{M}gf$. We denote the middle maps by $s'_{G_{M}f}:G_{M}X \rightarrow E'_{G_{M}f}$ 
     and $s'_{G_{M}gf}:G_{M}X \rightarrow E'_{G_{M}gf}$ .
     
    Now the whole diagram 
    \begin{equation} \label{eqn:octahed1}
       \begin{tikzcd}
        X \arrow{r}{i_{f}} \arrow{d}{\id_{X}}  & E_{f} \arrow{d}{s_{f}}  \\
        X \arrow{r}{i_{gf}} \arrow{d}{f} & E_{gf} \arrow{d}{s_{gf}} \\
        Y \arrow{r}{i_{g}} & E_{g}
       \end{tikzcd}
    \end{equation} 
      admits sections into 
      \begin{equation} \label{eqn:octahed2}
        \begin{tikzcd}
        F_{M}G_{M}X \arrow{r}{F_{M}G_{M}i_{f}} \arrow{d}{\id_{F_{M}G_{M}X}}  & F_{M}G_{M}E_{f} \arrow{d}{F_{M}G_{M}s_{f}}  \\
        F_{M}G_{M}X \arrow{r}{F_{M}G_{M}i_{gf}} \arrow{d}{F_{M}G_{M}f} & F_{M}G_{M}E_{gf} \arrow{d}{F_{M}G_{M}s_{gf}} \\
        F_{M}G_{M}Y \arrow{r}{F_{M}G_{M}{i_{g}}} & F_{M}G_{M}E_{g}
     \end{tikzcd}
      \end{equation} 
      by naturality of the sections, thereby extending into a direct summand of ocathedron associated to \eqref{eqn:octahed1} into oactahedron associated to 
      \eqref{eqn:octahed2}. Furthermore since the pairs $(\id_{F_{M}G_{M}X}, F_{M}t_{f})$, $(\id_{F_{M}G_{M}X}, F_{M}t_{gf})$ 
      and $(\id_{F_{M}G_{M}Y},  F_{M}t_{g})$ all satisfy hypothesis of corollary \ref{cor:indcisom} the octahedron associated to \eqref{eqn:octahed2} is isomorphic to the octahedron associated to \eqref{eqn:octahed3}. So the octahedron associated to \eqref{eqn:octahed1} becomes a direct summand of octahedron associated to \eqref{eqn:octahed3}.
    \begin{equation} \label{eqn:octahed3}
     \begin{tikzcd}
        F_{M}G_{M}X \arrow{r}{F_{M}i'_{G_{M}f}} \arrow{d}{\id_{F_{M}G_{M}X}}  & F_{M}E'_{G_{M}f} \arrow{d}{F_{M}s'_{G_{M}f}}  \\
        F_{M}G_{M}X \arrow{r}{F_{M}i'_{G_{M}gf}} \arrow{d}{F_{M}G_{M}f} & F_{M}E'_{G_{M}gf} \arrow{d}{F_{M}s'_{G_{M}gf}} \\
        F_{M}G_{M}Y \arrow{r}{F_{M}i'_{G_{M}g}} & F_{M}E'_{G_{M}g}
     \end{tikzcd}
    \end{equation} .
\end{remark}

We summarize the compuation done in the remark above as the following theorem.

\begin{theorem} \label{the:mainsepm}
Let $\cC$ be a cofibrantly generated stable model category, and let $M$ be a 
separable monad on $\cC$ such that $M$ satisfies the assumptions of theorem 
\ref{the:stabcomp}. Then $\module{M}$ is a cofibrantly generated stable model 
category. Moreover, a cofiber sequence admits a direct summand inclusion into a suitable
cofiber sequence.
\end{theorem}
\begin{proof} For the moreover part, we note that if $[f]:X  \rightarrow Y$ is a map in $\Ho(\module{M})$. We can represent it by some $\Tilde{f}: \Tilde{X} \rightarrow \Tilde{Y}$ in $\module{M}$, where $\Tilde{X}$ and $\Tilde{Y}$ are bifibrant objects in $\module{M}$. We conclude by lemma \ref{lem:dirincof} and remark \ref{rem:sepcompu}.
\end{proof}
\begin{remark} \label{rem:sepscmnd}
    \begin{enumerate}
        \item With the availability of sections along with stable compatibility, we note that to show a cofiber sequence  admits a direct summand inclusion as above, we did not appeal to the trick of lifting idempotents up to nilpotence as in \cite[lemma 1.6(c)]{balmer:septrcat}; yet we obtain a similar looking conclusion about the data of cofiber sequences-distinguished triangles as in \cite[theorem 4.1]{balmer:septrcat}.
        \item  By induction, we can now extend the same observation as above to an $n$-triangle (see \cite[definition 5.1]{balmer:septrcat}) for $n \ge 1$ as noted in \cite[remark 5.14]{balmer:septrcat} and again obtain similar looking conclusion as in \cite[theroem 5.17]{balmer:septrcat}.
    \end{enumerate}
\end{remark}
\section{Applications} \label{sec:applications}
For all the applications in this paper we are in the situation where a functor $G: \cD \rightarrow \cC$ admits both left and right adjoints i.e. $(F \dashv G \dashv H)$, this gives us two monads: $M=GF \colon \cC \rightarrow \cC$ and  $N=HG \colon \cD \rightarrow \cD$. Let $\module{M}$ and $\module{N}$ be the respective module categories. By the adjoint lifting theorem, we get $(F' \dashv G' \dashv H')$ the induced adjoints for $G': \module{N} \rightarrow \module{M}$. 

\subsection{Category of modules over related monads} \label{ssc:drumhack}

Here, we recall and improve upon a result of Drummond-Cole and Hackney.  We describe the similarities of properties shared by $\module{M}$ and $\module{N}$.

    We recall (Theorem 2.3 \cite{MR3937590}): Suppose $(\mathcal{M},\mathcal{W},I,J)$ is a cofibrantly generated model category and $\mathcal{N}$ is a bicomplete category. Let $F: \mathcal{N} \rightarrow \mathcal{M}$ be a functor which is right adjoint to $L$ and left adjoint to $R$. If $(FL,FR)$ is a Quillen adjunction, then there exists a right-induced cofibrantly generated model category structure $(\mathcal{N},F^{-1}\mathcal{W},LI,LJ)$. Both $(L,F)$ and $(F,R)$ are Quillen adjunctions.
    

    \begin{remark} \label{rem:autostcy}
    We point out that stable compatibility is always satisfied when the functor that creates the right induced structure (for example, for $G_{M}$ in section \ref{sec:stablecpt} in page \pageref{itm:rtmodstr}, and $F$ in theorem stated above), and is both a left and a right Quillen functor. So if $\mathcal{M}$ is stable then $\mathcal{N}$ is stable as well.
    \end{remark}
    
Let  $A$ be a cofibrant highly structured ring spectrum ($S$-module)
as defined in \cite{MR1417719} and $f \colon S \rightarrow A$ be the unit map. 
Denote by 
$\mathcal{M}_{S}$ the cofibrantly generated topological (and hence simplicial), stable model category of $S$-modules and 
let $\mathcal{M}_{A}$ be the cofibrantly generated topological (and hence 
simplicial), stable model category of $A$-modules.  The unit map induces 
$(f_{*}=A \wedge_{S}-) \dashv f^{*} \dashv F_{S}(A,-))$.

The functor $f^{*}$ is both a left and right Quillen functor. For $M=f^{*}f_{*}$, $
\module{M} = \mathcal{M}_{A}$ and so $F_{M}=f_{*}$ and $G_{M}=f^{*}$. We point out 
here that if $G_{M}$ is both a left and right Quillen functor, then it automatically 
satisfies our condition of stable compatibility.

So $\module{M}$ is always a cofibrantly generated simplicial and stable model category for any $A$. But, can we say anything about the category $\module{N}$?
So, if $(F' \dashv G' \dashv H')$ induced by $(f_{*}=A \wedge_{S}-) \dashv f^{*} \dashv F_{S}(A,-))$  is such that the triple satisfies \cite[theorem 2.3]{MR3937590}; and hence by remark \ref{rem:autostcy}, $\module{N}$ is also a cofibrantly generated simplicial and stable model category.

\subsection{Finite, flat maps between schemes} \label{ssc:finflmor}
 We now consider the monads in Algebraic geometry. We state the example for Noetherian schemes, but the results hold in greater generality.
 
 Let $f: X \rightarrow Y$ be a finite, flat morphism of Noetherian schemes, for such a suitable $f$ we get an adjunction 
 $\adjunction{f^{*}}{\Ch(QCoh(Y))}{Ch(QCoh(X))}{f_{*}}$ and the monad $M_{f}=f_{*}f^{*}$ on $Ch(QCoh(Y))$.
We consider the locally-free model structure on $\chA$ as stated in theorem \ref{thm:locallyfreeModelStruct} for $\mathscr{A} = \QCoh(Y)$.
We claim that when $f$ is finite and flat, it satisfies the assumptions of lemma \ref{lem:cndforsc}. To see this, we note that we have the adjunctions $(f^{*} \dashv f_{*} \dashv f^{!} )$, and therefore $M_{f}=f_{*}f^{*}$ preserves colimits. The pull back $f^*$ is known to preserve locally free sheaves, and so does $f_*$ for finite (hence affine) and flat $f$ (see, for example, 
\cite[\href{https://stacks.math.columbia.edu/tag/02K9}{Tag 02K9}, lemma 29.49.2]{stacks-project}; and also \cite[II, exercise 5.8; III, proposition 9.2]{hartshorne:alggeom} for the case of reduced, connected noetherian schemes).
Hence, $M_f$ takes locally free sheaves to locally free sheaves. The functor $M_f$ also preserve
subobjects as both $f^*$ and $f_*$ are exact under the assumptions.
Therefore, $M_f$ preserves the set of generating cofibrations and generating acyclic 
cofibrations as mentioned in \ref{def: I and J}. Further, by \cite[exercise III, 8.2]{hartshorne:alggeom}, we know that 
for affine $f$, $R^i f_* = 0$ for $i > 0$, and hence $R f_* = f_*$

Thus, the monad $M_f = f_* f^*$ on $\chA$ is
computed as $M_f(A_{\bullet}) = (\dotsc \rightarrow f_* f^*(A_i) \rightarrow f_* f^* 
(A_{i+1}) \rightarrow \dotsc)$. From this, it is clear that elements of $\module[\chA]
{M_f}$ pairs $(A_{\bullet}, \lambda_{\bullet})$ with $\lambda_n \colon f_* f^* A_n 
\longrightarrow A_n$, where all the squares in the corresponding diagram commute.
The above description tells us that $\module[\chA]{M_f}$ is the same as 
$\Ch(\module[\cA]{M_f})$.  
Clearly, in this case, $T$ is the shift functor and hence it commutes with the
forgetful functor $G_M$. 
By theorem \ref{the:stabcomp}, we get that $\module[\chA]{M_f}$ is a 
cofibrantly generated stable model category.

Examples of such maps are \'etale maps \cite{balmer:etalneth} and the 
$n$-power map $z \mapsto z^n$.

\subsection{A remark about comonads and comodules over them} \label{ssc:comodules}

By considering monads in the opposite category, one gets the notion of 
comonads, and one can construct comodules over them such that the opposite 
category of comodules over a comonad $W$ over $\cC$ coincide with the module 
category of $W^{\text{op}}$ over $\cC^{\text{op}}$ 
(compare \cite[remark 5.19]{balmer:septrcat}). 
Thus, by considering the corresponding notion of stable compatibility, namely
that in the situation of a left-transferred model structure the suspension should commute with the forgetful functor from the 
category of comodules, and that the forgetful functor should preserve fibrant
objects, we would get a criterion for the category of comodules to be stable.

\bibliography{references}
\bibliographystyle{amsalpha}


\end{document}
